% \section*{Essential Implications and Proofs}

{\scriptsize
\renewcommand{\arraystretch}{1}
\setlength{\tabcolsep}{2pt}

\noindent
% ============================================================
% LEFT MINIPAGE
% ============================================================
\begin{minipage}[t]{0.495\textwidth}
\vspace{0pt}

\textbf{\large Essential Implications}\par
\vspace{2pt}

\begin{tabular}{
    @{}
    >{\bfseries\raggedright\arraybackslash}p{2.9cm}
    >{\raggedright\arraybackslash}p{10.5cm}
    @{}
}
\toprule

Regularity &
$C^\infty \Rightarrow C^n \Rightarrow C^1
\Rightarrow \text{differentiable}
\Rightarrow \text{continuous}$ \\

Derivative bound &
$|f'|\le M
\Rightarrow \text{Lipschitz}
\Rightarrow \text{uniformly continuous}
\Rightarrow \text{continuous}$ \\

Derivative sign &
$f'>0 \Rightarrow f\text{ strictly increasing}
\Rightarrow f\text{ injective}$,
\quad
$f'<0 \Rightarrow f\text{ strictly decreasing}
\Rightarrow f\text{ injective}$ \\

Monotonicity &
$f'\ge0 \Rightarrow f\text{ increasing}$,
\quad
$f'\le0 \Rightarrow f\text{ decreasing}$ \\

Constant function &
$f'=0 \Longleftrightarrow f=\text{constant}$,
\qquad
$f'=g' \Rightarrow f-g=C$ \\

Inverse &
$f\text{ strictly monotone}
\Rightarrow f\text{ injective}
\Rightarrow f^{-1}\text{ exists on }f(D)$ \\

Continuous inverse &
$f:I\to J\text{ continuous and bijective}
\Rightarrow f^{-1}:J\to I\text{ continuous}$ \\

Extremum condition &
$x_0\text{ interior local extremum}+
f\text{ differentiable at }x_0
\Rightarrow f'(x_0)=0$ \\

First derivative test &
$f':-\to+ \Rightarrow x_0\text{ local minimum}$,
\qquad
$f':+\to- \Rightarrow x_0\text{ local maximum}$ \\

Second derivative test &
$f'(x_0)=0,\ f''(x_0)>0 \Rightarrow x_0\text{ local minimum}$,
\quad
$f'(x_0)=0,\ f''(x_0)<0 \Rightarrow x_0\text{ local maximum}$ \\

Convexity &
$f''\ge0 \Rightarrow f'\text{ increasing}\Rightarrow f\text{ convex}$,
\qquad
$f''\le0 \Rightarrow f\text{ concave}$ \\

Convex minimum &
$f\text{ convex},\ f'(x_0)=0
\Rightarrow x_0\text{ global minimum}$ \\

Continuous compact &
$f\in C([a,b])
\Rightarrow f\text{ bounded and attains }\min,\max$ \\

Intermediate value &
$f\in C([a,b]),\ f(a)f(b)<0
\Rightarrow \exists c\in(a,b):f(c)=0$ \\

Unique root &
$\text{continuity}+\text{sign change}+\text{strict monotonicity}
\Rightarrow \exists!c:f(c)=0$ \\

Sequential continuity &
$f\text{ continuous at }x_0
\Longleftrightarrow
\bigl(x_n\to x_0\Rightarrow f(x_n)\to f(x_0)\bigr)$ \\

Sequence convergence &
$a_n\to L
\Rightarrow (a_n)\text{ bounded}$ \\

Subsequences &
$a_n\to L
\Rightarrow a_{n_k}\to L
\text{ for every subsequence}$ \\

Divergence via subsequences &
$a_{n_k}\to A,\ a_{m_k}\to B,\ A\neq B
\Rightarrow (a_n)\text{ diverges}$ \\

Accumulation point &
$A\text{ accumulation point}
\Longleftrightarrow
\exists(a_{n_k}):a_{n_k}\to A$ \\

Bolzano--Weierstrass &
$(a_n)\text{ bounded}
\Rightarrow \exists\text{ convergent subsequence}$ \\

Monotone convergence &
$\text{monotone}+\text{bounded}
\Rightarrow \text{convergent}$ \\

Increasing sequence &
$a_n\uparrow,\ (a_n)\text{ bounded above}
\Rightarrow a_n\to\sup\{a_n:n\in\mathbb N\}$ \\

Decreasing sequence &
$a_n\downarrow,\ (a_n)\text{ bounded below}
\Rightarrow a_n\to\inf\{a_n:n\in\mathbb N\}$ \\

Cauchy &
$\text{in }\mathbb R:\quad
(a_n)\text{ convergent}
\Longleftrightarrow
(a_n)\text{ Cauchy}
\Rightarrow
(a_n)\text{ bounded}$ \\

Order and limits &
$a_n\le b_n\text{ eventually},\
a_n\to A,\ b_n\to B
\Rightarrow A\le B$ \\

Sandwich theorem &
$a_n\le b_n\le c_n,\ a_n\to L,\ c_n\to L
\Rightarrow b_n\to L$ \\

Useful squeeze form &
$|b_n|\le c_n,\ c_n\to0
\Rightarrow b_n\to0$ \\

Limsup/liminf &
$(a_n)\text{ bounded}:\quad
a_n\text{ converges}
\Longleftrightarrow
\liminf a_n=\limsup a_n$ \\

% \bottomrule

Series term test &
$\sum a_n\text{ converges}
\Rightarrow a_n\to0$ \\

Series divergence &
$a_n\not\to0
\Rightarrow \sum a_n\text{ diverges}$ \\

Absolute convergence &
$\sum |a_n|\text{ converges}
\Rightarrow
\sum a_n\text{ converges}
\Rightarrow
a_n\to0$ \\

Comparison test &
$0\le b_n\le a_n,\ \sum a_n\text{ converges}
\Rightarrow \sum b_n\text{ converges}$ \\

Divergence comparison &
$0\le b_n\le a_n,\ \sum b_n\text{ diverges}
\Rightarrow \sum a_n\text{ diverges}$ \\

Ratio test &
$\lim\left|\frac{a_{n+1}}{a_n}\right|<1
\Rightarrow \sum a_n\text{ absolutely convergent}$,
\quad
$>1\Rightarrow\text{divergent}$ \\

Root test &
$\limsup |a_n|^{1/n}<1
\Rightarrow \sum a_n\text{ absolutely convergent}$,
\quad
$>1\Rightarrow\text{divergent}$ \\

Power series &
$|x-a|<R\Rightarrow\text{absolute convergence}$,
\quad
$|x-a|>R\Rightarrow\text{divergence}$,
\quad
$|x-a|=R\Rightarrow\text{check separately}$ \\

Inside radius &
$|x-a|<R
\Rightarrow \text{termwise differentiation/integration}
\Rightarrow C^\infty$ \\

Riemann integrability &
$f\text{ continuous on }[a,b]
\Rightarrow f\text{ Riemann integrable}
\Rightarrow f\text{ bounded}$ \\

Monotone integrability &
$f\text{ monotone on }[a,b]
\Rightarrow f\text{ Riemann integrable}$ \\

Integral order &
$f\le g
\Rightarrow
\int_a^b f(x)\,dx
\le
\int_a^b g(x)\,dx$ \\

Integral bound &
$|f|\le M
\Rightarrow
\left|\int_a^b f(x)\,dx\right|
\le M|b-a|$ \\

Antiderivatives &
$F'=f,\ G'=f
\Rightarrow F-G=C$ \\

FTC I &
$f\in C([a,b]),\
F(x)=\int_a^x f(t)\,dt
\Rightarrow F'=f$ \\

FTC II &
$F'=f
\Rightarrow
\int_a^b f(x)\,dx=F(b)-F(a)$ \\

Uniform convergence &
$f_n\xrightarrow{\mathrm{unif}}f
\Rightarrow
f_n\xrightarrow{\mathrm{pointwise}}f$ \\

Uniform + continuity &
$f_n\text{ continuous},\
f_n\xrightarrow{\mathrm{unif}}f
\Rightarrow f\text{ continuous}$ \\

Uniform + integration &
$f_n\text{ integrable},\
f_n\xrightarrow{\mathrm{unif}}f
\Rightarrow
\displaystyle
\lim_{n\to\infty}\int_a^b f_n
=
\int_a^b f$ \\

\midrule

False converse &
$\text{continuous}\not\Rightarrow\text{differentiable}\ (|x|)$ \\

False converse &
$\text{bounded sequence}\not\Rightarrow\text{convergent}\ ((-1)^n)$ \\

False converse &
$a_n\to0\not\Rightarrow\sum a_n\text{ convergent}\ 
\left(\sum\frac1n\right)$ \\

False converse &
$f'(x_0)=0\not\Rightarrow x_0\text{ extremum}\ (x^3)$ \\

False converse &
$f\text{ strictly increasing}\not\Rightarrow f'>0\ (x^3)$ \\

False converse &
$\text{pointwise}\not\Rightarrow\text{uniform}$ \\

False converse &
$\text{Riemann integrable}\not\Rightarrow\text{continuous}$ \\

False converse &
$\text{convergent series}\not\Rightarrow
\text{absolutely convergent}$ \\

Regularity &
$C^\infty\Rightarrow C^n\Rightarrow C^1
\Rightarrow\text{differentiable}
\Rightarrow\text{continuous}$ \\

% Paste all remaining implication rows here.

\bottomrule

\addlinespace[2mm]

\bottomrule

Continuity of composition (proof) &
Let $f$ be continuous at $x_0$ and $g$ continuous at $f(x_0)$. Given
$\varepsilon>0$, choose $\eta>0$ such that
$
|y-f(x_0)|<\eta
\Rightarrow
|g(y)-g(f(x_0))|<\varepsilon.
$
Choose $\delta>0$ such that
$
|x-x_0|<\delta
\Rightarrow
|f(x)-f(x_0)|<\eta.
$
Then
$
|x-x_0|<\delta
\Rightarrow
|g(f(x))-g(f(x_0))|<\varepsilon.
$
\\
\bottomrule
\end{tabular}

\end{minipage}
\hfill
% ============================================================
% RIGHT MINIPAGE
% ============================================================
\begin{minipage}[t]{0.495\textwidth}
\vspace{0pt}

\textbf{\large Essential Proofs I}\par
\vspace{2pt}

\begin{tabular}{
    @{}
    >{\bfseries\raggedright\arraybackslash}p{3.0cm}
    >{\raggedright\arraybackslash}p{10.4cm}
    @{}
}
\toprule

Uniqueness of the limit &
Assume $a_n\to A$ and $a_n\to B$. For arbitrary $\varepsilon>0$, choose
$N_A,N_B$ such that
$n\ge N_A\Rightarrow |a_n-A|<\varepsilon/2$ and
$n\ge N_B\Rightarrow |a_n-B|<\varepsilon/2$.
For $n\ge N:=\max\{N_A,N_B\}$,
$
|A-B|
\le |A-a_n|+|a_n-B|
<\varepsilon.
$
Since this holds for every $\varepsilon>0$, $|A-B|=0$, hence $A=B$.
\\

\midrule

Subsequences preserve limits &
Assume $a_n\to L$ and let $(a_{n_k})$ be a subsequence. For arbitrary
$\varepsilon>0$, choose $N$ such that
$
n\ge N\Rightarrow |a_n-L|<\varepsilon.
$
Since $(n_k)$ is strictly increasing, $n_k\ge k$. Thus, for $k\ge N$,
$
n_k\ge k\ge N
\Rightarrow
|a_{n_k}-L|<\varepsilon.
$
Therefore $a_{n_k}\to L$.
\\

\midrule

Accumulation point gives infinitely many terms &
Let $A$ be an accumulation point. Then there exists a subsequence
$a_{n_k}\to A$. For arbitrary $\varepsilon>0$, choose $K$ such that
$
k\ge K\Rightarrow |a_{n_k}-A|<\varepsilon.
$
Hence
$
a_{n_k}\in(A-\varepsilon,A+\varepsilon)
$
for every $k\ge K$. Since the indices $n_k$ are distinct and there are
infinitely many $k\ge K$, the interval contains infinitely many sequence
terms.
\\

\midrule

Convergent sequences are bounded &
Assume $a_n\to L$. Choose $N$ such that
$
n\ge N\Rightarrow |a_n-L|<1.
$
For $n\ge N$,
$
|a_n|
\le |a_n-L|+|L|
<1+|L|.
$
Define
$
M:=\max\{|a_0|,\ldots,|a_{N-1}|,1+|L|\}.
$
Then $|a_n|\le M$ for every $n\in\mathbb N_0$; hence $(a_n)$ is bounded.
\\

\midrule

Sandwich theorem &
Assume $a_n\to L$, $c_n\to L$, and
$a_n\le b_n\le c_n$ for all $n\ge N_0$.
For arbitrary $\varepsilon>0$, choose $N_1,N_2$ such that
$
n\ge N_1\Rightarrow L-\varepsilon<a_n<L+\varepsilon
$
and
$
n\ge N_2\Rightarrow L-\varepsilon<c_n<L+\varepsilon.
$
For
$
n\ge N:=\max\{N_0,N_1,N_2\},
$
$
L-\varepsilon<a_n\le b_n\le c_n<L+\varepsilon.
$
Thus $|b_n-L|<\varepsilon$, so $b_n\to L$.
\\

\midrule

Useful squeeze form &
If $|b_n|\le c_n$ eventually and $c_n\to0$, then
$
-c_n\le b_n\le c_n.
$
Since $-c_n\to0$ and $c_n\to0$, the sandwich theorem gives $b_n\to0$.
\\

\midrule

Order preservation under limits &
Assume $a_n\le b_n$ eventually, $a_n\to A$, and $b_n\to B$.
Suppose $A>B$ and set $\varepsilon:=(A-B)/3>0$.
Eventually,
$
a_n>A-\varepsilon
$
and
$
b_n<B+\varepsilon.
$
But
$
A-\varepsilon>B+\varepsilon,
$
so eventually $a_n>b_n$, contradicting $a_n\le b_n$. Hence $A\le B$.
\\

\midrule

Unique accumulation point of a convergent sequence &
If $a_n\to L$, every convergent subsequence also converges to $L$.
Since every accumulation point is the limit of some subsequence, every
accumulation point equals $L$. Moreover, $(a_n)$ itself converges to $L$,
so $L$ is an accumulation point. Hence $L$ is the unique accumulation
point.
\\

\midrule

Nonnegative series &
Let
$
s_n:=\sum_{k=0}^{n}a_k
$
with $a_k\ge0$. Then
$
s_{n+1}-s_n=a_{n+1}\ge0,
$
so $(s_n)$ is increasing. Hence
$
(s_n)\text{ bounded}
\Rightarrow
(s_n)\text{ convergent}
\Rightarrow
\sum_{n=0}^{\infty}a_n\text{ converges}.
$
If $(s_n)$ is unbounded, then the series diverges.
\\

\midrule

Necessary condition for series convergence &
Assume
$
\sum_{n=0}^{\infty}a_n
$
converges and let
$
s_n:=\sum_{k=0}^{n}a_k\to S.
$
Since
$
a_n=s_n-s_{n-1},
$
we obtain
$
\lim_{n\to\infty}a_n
=
\lim_{n\to\infty}s_n
-
\lim_{n\to\infty}s_{n-1}
=
S-S=0.
$
Thus
$
\sum a_n\text{ convergent}
\Rightarrow
a_n\to0.
$
\\

\midrule

Cauchy criterion: convergence $\Rightarrow$ Cauchy &
Assume $a_n\to L$. For arbitrary $\varepsilon>0$, choose $N$ such that
$
n\ge N\Rightarrow |a_n-L|<\varepsilon/2.
$
Then for all $m,n\ge N$,
$
|a_n-a_m|
\le
|a_n-L|+|a_m-L|
<
\varepsilon.
$
Hence $(a_n)$ is Cauchy.
\\

\midrule

Cauchy criterion: Cauchy $\Rightarrow$ convergence &
Let $(a_n)$ be Cauchy. Then $(a_n)$ is bounded, so by
Bolzano--Weierstrass there exists a subsequence
$
a_{n_k}\to L.
$
For $\varepsilon>0$, choose $N_1$ such that
$
m,n\ge N_1\Rightarrow |a_n-a_m|<\varepsilon/2,
$
and choose $K$ such that
$
k\ge K\Rightarrow |a_{n_k}-L|<\varepsilon/2.
$
Choose $k$ with
$
k\ge K
$
and
$
n_k\ge N_1.
$
Then for every $n\ge N_1$,
$
|a_n-L|
\le
|a_n-a_{n_k}|+|a_{n_k}-L|
<
\varepsilon.
$
Thus $a_n\to L$.
\\

\midrule

Cauchy sequences are bounded &
Let $(a_n)$ be Cauchy. Choose $N$ such that
$
m,n\ge N\Rightarrow |a_n-a_m|<1.
$
Fix $m=N$. Then for $n\ge N$,
$
|a_n|
\le
|a_n-a_N|+|a_N|
<
1+|a_N|.
$
Set
$
M:=\max\{|a_0|,\ldots,|a_{N-1}|,1+|a_N|\}.
$
Then
$
|a_n|\le M
$
for every $n$, so $(a_n)$ is bounded.
\\

\midrule

Monotone sequence with bounded subsequence &
Assume $(a_n)$ is increasing and has a bounded subsequence
$
|a_{n_k}|\le M.
$
For every $n$, choose $k$ with $n\le n_k$. By monotonicity,
$
a_n\le a_{n_k}\le M.
$
Thus $(a_n)$ is bounded above. Since it is increasing and bounded, it
converges. The decreasing case is analogous.
\\

\midrule

Differentiable $\Rightarrow$ continuous &
Assume $f$ is differentiable at $x_0$. For $h\neq0$,
$
f(x_0+h)-f(x_0)
=
\frac{f(x_0+h)-f(x_0)}{h}\,h.
$
Hence
$
\lim_{h\to0}\bigl(f(x_0+h)-f(x_0)\bigr)
=
f'(x_0)\cdot0=0.
$
Thus
$
\lim_{x\to x_0}f(x)=f(x_0),
$
so $f$ is continuous at $x_0$.
\\

\midrule

Rolle's theorem &
Since $f\in C([a,b])$, it attains a minimum and maximum. If $f$ is
constant, then $f'=0$ on $(a,b)$. Otherwise, because $f(a)=f(b)$, at
least one extremum occurs at some $\xi\in(a,b)$. Since $f$ is
differentiable there, Fermat's theorem gives
$
f'(\xi)=0.
$
\\

\midrule

Mean value theorem &
Define
$
g(x):=f(x)-\frac{f(b)-f(a)}{b-a}(x-a).
$
Then $g$ is continuous on $[a,b]$, differentiable on $(a,b)$, and
$
g(a)=g(b)=f(a).
$
By Rolle, some $\xi\in(a,b)$ satisfies
$
0=g'(\xi)
=
f'(\xi)-\frac{f(b)-f(a)}{b-a}.
$
Hence
$
f'(\xi)=\frac{f(b)-f(a)}{b-a}.
$
\\

\midrule

Strict monotonicity $\Rightarrow$ injectivity &
Let $x_1\neq x_2$; w.l.o.g. $x_1<x_2$. If $f$ is strictly increasing,
then $f(x_1)<f(x_2)$; if strictly decreasing, then
$f(x_1)>f(x_2)$. Thus $f(x_1)\neq f(x_2)$, so $f$ is injective.
\\

\midrule

Algebra of continuous functions &
Let $x_n\to x_0$. Continuity gives
$
f(x_n)\to f(x_0)
$
and
$
g(x_n)\to g(x_0).
$
By the limit laws,
$
f(x_n)\pm g(x_n)\to f(x_0)\pm g(x_0),
$
$
f(x_n)g(x_n)\to f(x_0)g(x_0),
$
and, if $g(x_0)\neq0$,
$
f(x_n)/g(x_n)\to f(x_0)/g(x_0).
$
Hence the corresponding functions are continuous.
\\

\bottomrule
\end{tabular}

\end{minipage}
}